\documentclass[10pt,a4paper]{amsart}
\usepackage[margin=19mm]{geometry}
\usepackage{amsmath,amssymb,mathtools}
\usepackage[scr=rsfs,cal=euler]{mathalfa}
\usepackage{xfrac}
\usepackage[unicode,hidelinks,bookmarks=false]{hyperref}
\newcommand{\ie}{i.e.\ }
\newcommand{\cf}{cf.\ }
\newcommand{\resp}{respectively\ }
\newcommand{\Subsection}{\subsection}
\providecommand{\nobreakdash}{\nobreak\hbox{-}}
\setlength{\emergencystretch}{2em}
\multlinegap0pt
\usepackage{bm}
\usepackage{bbm}
\usepackage{dsfont}
\usepackage{xspace}
\usepackage{cancel}
\usepackage{mathtools}
\usepackage{amsthm}
\makeatletter
\@ifundefined{theorem}{\newtheorem{theorem}{Theorem}}{}
\@ifundefined{lemma}{\newtheorem{lemma}[theorem]{Lemma}}{}
\@ifundefined{proposition}{\newtheorem{proposition}[theorem]{Proposition}}{}
\makeatother
\title[On comass and stable systolic inequalities]{On comass and stable systolic inequalities}
\author{James J. Hebda, Mikhail G. Katz}
\date{Source: arXiv:2411.13966v2 (CC BY 4.0)}
\begin{document}
\maketitle

\noindent\emph{Source and licence.} On comass and stable systolic inequalities. Version arXiv:2411.13966v2 (CC BY 4.0), distributed under the Creative Commons Attribution 4.0 International licence, as recorded in the arXiv licence metadata for this exact version and shown on the abstract page https://arxiv.org/abs/2411.13966v2. Full rights notice: licenses/arxiv-2411.13966-CC-BY-4.0.txt.\par\smallskip

\noindent\textit{Editorial scope.} Counted unit: the Proposition whose source label is \texttt{p:i2} in the article (source0.tex lines 190--195, source label \texttt{p:i2}), delivered together with the article's own complete original proof of it (source0.tex lines 197--252, one \texttt{proof} environment of 56 lines). No mathematics is written, repaired or re-derived: statement and proof are the article's own text line for line. Required context retained uncounted: none. Every definition, display and label of those lines is retained unchanged. Omitted from this package and not redistributed: 1--189, 196--196, 253--542. Nothing inside a retained excerpt is omitted or altered: every retained body is delivered exactly as the article's lines are written, so no omission receipt of any kind accompanies this package. Citations: the retained excerpts carry no citation command at all, so no bibliography entry of the article is transcribed and no reference is redistributed. Reference adaptation: none was needed and none was made; every reference inside the delivered unit resolves inside it, so no unresolved reference or citation is produced. Printed slips kept literal: no printed word, symbol, hypothesis or number of the counted statement or of its proof was altered. Layout and class: the article's own class is \texttt{amsart} with packages including amsmath, amsthm, geometry, graphicx, amssymb, epstopdf, xcolor, hyperref, breakurl; the wrapper is the lane's pinned \texttt{amsart} editorial wrapper at 10pt on A4, which restates the theorem-like environments the retained text needs and adds the article's own macro definitions that the retained text needs (none), with extra wrapper packages (bm, bbm, dsfont, xspace, cancel, mathtools). The delivered numbering is the unit's own. Assets: no retained span contains a figure environment, a graphics-inclusion command, a PDF-page inclusion, a table or a TikZ picture; no image file is copied into this package, the package declares no asset dependency and no image file is redistributed with it. Licence: version arXiv:2411.13966v2 of this article is distributed under the Creative Commons Attribution 4.0 International licence, as recorded in the arXiv licence metadata for this exact version and shown on the abstract page https://arxiv.org/abs/2411.13966v2. A full rights notice is delivered at licenses/arxiv-2411.13966-CC-BY-4.0.txt. Characters: every retained excerpt is pure ASCII, so the delivered engine, XeLaTeX with its default Latin Modern text font, reports no missing character.
\begin{proposition}\label{p:i2}%2.5
Let $1\leq k<p<n$, then
\begin{equation*}
C_{n,p}^2 \leq \frac {\binom{n}{p-k}}{\binom{p}{k}} C_{n,k}^2.
\end{equation*}
\end{proposition}

\begin{proof}
Let $\phi = \sum_{I} \beta_I e_I^\ast  \in \Lambda^p(\mathbf{R}^n)^\ast$ where $I$ runs over all multi--indices of length $p$, that is, $I$ is a strictly increasing sequence of $p$ integers between $1$ and $n$.

Fix a multi--index $J$ of length $p-k$ of integers between $1$ and
$n$. Suppose $I$ is a multi--index of length $p$ that contains $J$ as
a subsequence.  In this case write $I \supset J$. Let $I\backslash J$
denote the multi--index of length $k$ complementary to $J$ in
$I$. Thus $e_{I\backslash J} \wedge e_J = \pm e_I \in
\Lambda^p(\mathbf{R}^n)$ where the sign depends on the permutation of
the elements in $I$. Pull out the terms of $\phi$ that contain all the
indices in $J$ by setting
\begin{equation*}
\phi_J  =  \sum_{I\supset J}  \beta_I e_I^\ast.
\end{equation*}
Then let
\begin{equation*}
\psi_J = \sum_{I\supset J} \pm \beta_I e^\ast_{I\backslash J} \in \Lambda^k(\mathbf{R}^n)^\ast,
\end{equation*}
 where  the sign is chosen to agree with the sign in $e_{I\backslash J} \wedge e_J = \pm e_I$, and the sums run over all multi-indices $I$ of length $p$ with $I \supset J$.
 Thus $\phi_J = \psi_J\wedge e_J^\ast$ and $ | \psi_J|^\ast =  \sqrt{\sum_{I \supset J} \beta_I^2}$.
 
 \begin{lemma}%2.6
 \begin{equation}\label{e:i2}
   \Vert \phi\Vert^{\ast2} C_{n,k}^2 \geq\sum_{I \supset J} \beta_I^2.
  \end{equation}
 \end{lemma}
 
 \begin{proof}
 Choose a simple $k$--vector  $\xi_J \in \Lambda^k(\mathbf{R}^n)$ such that $|\xi_J|=1$ and $ \psi_J(\xi_J) = \Vert \psi_J\Vert^\ast$. Let $V = \mathrm{span}\{ e_i : i \notin J\} \subset \mathbf{R}^n$. Since by construction $ \psi_J \in \Lambda^k(V)^\ast$, we may assume that $ \xi_J \in \Lambda^k(V) \subset \Lambda^k(\mathbf{R}^n)$.  Since $e_J$ is a simple $p-k$ vector of norm $1$, $ \xi_J \wedge e_J$ is a simple $p$-vector with $| \xi_J \wedge e_J|=1$ . Noting that $e_I^\ast(\xi_J\wedge e_J)=0$ when $I\not\supset J$, it follows that
 \begin{equation*}\Vert \phi\Vert^\ast \geq  \phi( \xi_J \wedge e_J)=  \phi_J( \xi_J \wedge e_J) = (\psi_J\wedge e_J^\ast)( \xi_J \wedge e_J)=\psi_J(\xi_J)e_J^\ast(e_J)  =  \Vert \psi_J\Vert^\ast.
 \end{equation*}
Hence, by definition of $C_{n,k}$,
\begin{equation*}
  \sqrt{\sum_{I \supset J} \beta_I^2} = |\psi_J|^\ast \leq C_{n,k} \Vert \psi_J\Vert^\ast \leq C_{n,k} \Vert \phi \Vert^\ast.
 \end{equation*} 
On squaring this inequality,
 \begin{equation*}
 \quad  \Vert \phi\Vert^{\ast2} C_{n,k}^2 \geq \sum_{I \supset J} \beta_I^2,
 \end{equation*}
 which completes the proof of the lemma.
 \end{proof}
 
 Returning to the proof of Proposition \ref{p:i2}, summing the inequality (\ref{e:i2}) over all $J$ gives
 \begin{equation*}
  \binom{n}{p-k}  C_{n,k}^2 \Vert \phi\Vert^{\ast2} \geq \sum_J\sum_{I \supset J} \beta_I^2 = \binom{p}{k}\sum_I \beta_I^2 =\binom{p}{k} |\phi|^{\ast 2}
  \end{equation*} 
 since in the double sum each $\beta_I^2$ occurs $\binom{p}{k}$ times, because there are $\binom{p}{k}$ multi--indices $J$ contained in each $I$.
 Hence
 \begin{equation*}
  \frac {|\phi|^{\ast 2}}{\Vert \phi\Vert^{\ast2}} \leq   \frac {\binom{n}{p-k}}{\binom{p}{k}} C_{n,k}^2. 
  \end{equation*}
 Therefore, since $\phi$ was arbitrary,
 \begin{equation*}
 C_{n,p}^2 \leq \frac {\binom{n}{p-k}}{\binom{p}{k}} C_{n,k}^2.
 \end{equation*}
 \end{proof}

\end{document}
